\documentclass[12pt]{article}
\usepackage{amsmath,amssymb,siunitx,graphicx}
\newcommand{\vect}[1]{\mathbf{#1}}
\newcommand{\dd}{\mathrm{d}}
\begin{document}
\section{Maxwell's equations}
In vacuum, with $\vect{E}$ the electric and $\vect{B}$ the magnetic field:
\begin{align}
\nabla\cdot\vect{E} &= \frac{\rho}{\varepsilon_0}, \label{eq:gauss}\\
\nabla\times\vect{B} &= \mu_0\vect{J} + \mu_0\varepsilon_0\frac{\partial\vect{E}}{\partial t}. \label{eq:ampere}
\end{align}
Equation~\eqref{eq:gauss} is Gauss's law. A charge of \SI{1.6e-19}{\coulomb} produces at \SI{1}{\metre} a field of about \SI{1.4e-9}{\volt\per\metre}.
\subsection{Units}
\begin{tabular}{lll}
\hline
Quantity & Symbol & Unit \\
\hline
Charge & $q$ & C \\
Field & $E$ & V/m \\
\hline
\end{tabular}
\begin{figure}
\centering
\includegraphics[width=6cm]{field-lines}
\caption{Field lines of a point charge.}
\end{figure}
\end{document}
